% Imported q=4 cone and the finite-mean Galton--Watson arguments.
% Uses only standard LaTeX/amsmath/amssymb/amsthm commands.
% Required theorem environments: proposition, lemma, theorem, remark.
% Bibliography key: OpenAI4.

\section{Four states: a posterior cone}
\label{sec:four-state-cone}

We use the four-state posterior cone and global pair inequalities of \cite{OpenAI4} and derive their consequences for arbitrary finite-mean offspring. We state the inequalities and describe their finite certificates explicitly.

For a posterior $p\in\Delta_4$, let
\[
 \zeta=Lp,\qquad
 L=\begin{pmatrix}
  1&1&-1&-1\\
  1&-1&1&-1\\
  1&-1&-1&1
 \end{pmatrix}.
\]
Define
\begin{align}
 A&=\zeta_1^2+\zeta_2^2+\zeta_3^2
      =4\sum_{i=1}^4p_i^2-1,
 & B&=\zeta_1\zeta_2\zeta_3,
 & C&=\zeta_1^2\zeta_2^2+\zeta_1^2\zeta_3^2+
       \zeta_2^2\zeta_3^2.
 \label{eq:q4-basic-polynomials}
\end{align}
In particular $0\le A\le3$.  The three polynomials needed below are
\begin{align}
 H_4={}&B-\frac1{10}A^2-\frac3{20}C+\frac3{10}AB
       -\frac1{50}A^3+\frac7{100}AC-\frac{51}{100}B^2,
       \label{eq:q4-H}\\
 F_4={}&4A-3B+\frac7{10}A^2-4C-\frac45AB,
       \label{eq:q4-F}\\
 G_4={}&1-\frac{161}{100}A-\frac{31}{10}B
       +\frac{773}{1000}A^2+4C-\frac{199}{50}AB.
       \label{eq:q4-G}
\end{align}
We regard these both as functions of $p$ and as polynomials in $\zeta=Lp$;
the argument will indicate which representation is used.

\subsection{The four-state inequalities}

For $p,r\in\Delta_4$ and $\pi\in S_4$, write
\[
 U_i^\pi=4p_i r_{\pi(i)},\qquad
 z_\pi=\sum_{i=1}^4U_i^\pi,\qquad
 P_i^\pi=\frac{U_i^\pi}{z_\pi}\quad(z_\pi>0).
\]
For a bounded function $f$ on $\Delta_4$, use the notation
\[
 [zf(P)]_4(p,r)=\frac1{24}\sum_{\pi\in S_4}z_\pi f(P^\pi),
\]
where a summand with $z_\pi=0$ is defined to be zero.  Thus the brackets
average over all twenty-four relative color permutations.

\begin{proposition}[Four-state polynomial inequalities, \cite{OpenAI4}]
\label{prop:q4-imported-pair}
The polynomials \eqref{eq:q4-H}--\eqref{eq:q4-G} satisfy
\[
 H_4\ge-\frac12,\qquad F_4\ge0,\qquad G_4\ge0
 \quad\hbox{on }\Delta_4.
\]
For $0\le t\le1$ and $\zeta\in L\Delta_4$,
\begin{equation}
 H_4(t\zeta)\ge t^3H_4(\zeta).
 \label{eq:q4-radial}
\end{equation}
For every $p,r\in\Delta_4$,
\begin{align}
 [zA(P)]_4(p,r)
 &\le A(p)+A(r)-H_4(p)F_4(r)-F_4(p)H_4(r)\notag\\
 &\hspace{22mm}
 -\frac{A(p)A(r)\bigl(A(p)+A(r)\bigr)}{1000},
 \label{eq:q4-pair-A}\\
 [zH_4(P)]_4(p,r)
 &\ge H_4(p)G_4(r)+G_4(p)H_4(r).
 \label{eq:q4-pair-H}
\end{align}
\end{proposition}

\subsection{Finite certificates for the four-state inequalities}
\label{subsec:q4-certificate}

We summarize the exact verification in \cite{OpenAI4} to specify precisely
the four-state computational ingredient.  Permuting the four probabilities
acts on $Lp$ by a signed permutation of three coordinates, with the product
of the three signs equal to one.  Consequently $A,B,C$ and
$H_4,F_4,G_4$ are invariant.  The averaged pair expressions are invariant
under separate permutations of $p$ and $r$.

It therefore suffices to work on the ordered simplex
\[
 \mathcal S_4=\{p\in\Delta_4:p_1\ge p_2\ge p_3\ge p_4\},
\]
whose vertices are
\[
 e_1,\quad \frac{e_1+e_2}{2},\quad
 \frac{e_1+e_2+e_3}{3},\quad u_4.
\]
Indeed its barycentric coordinates are
$p_1-p_2,2(p_2-p_3),3(p_3-p_4),4p_4$.
On a simplex, a homogeneous polynomial with nonnegative ordinary monomial
coefficients is nonnegative.  Homogenization by powers of the sum of the
barycentric coordinates does not alter values on the simplex.  This gives
the following exact scalar certificates; the displayed minimum refers to
ordinary homogeneous coefficients.
\begin{center}
\begin{tabular}{lrrr}
\hline
Polynomial & Degree & Coefficients & Minimum\\
\hline
$H_4+1/2$ & 6 & 84 & $19/50$\\
$F_4$ & 5 & 56 & $0$\\
$G_4$ & 14 & 680 & $87/1000$\\
$(3-\mathcal D)H_4$ & 6 & 84 & $0$\\
\hline
\end{tabular}
\end{center}
Here $\mathcal D=\sum_{j=1}^3\zeta_j\partial_{\zeta_j}$ is the Euler
operator.  Its scalar certificate implies \eqref{eq:q4-radial}, because
\[
 \frac{d}{dt}\bigl(t^{-3}H_4(t\zeta)\bigr)
 =t^{-4}(\mathcal D-3)H_4(t\zeta)\le0
 \qquad(0<t\le1),
\]
and $H_4(0)=0$ handles $t=0$.

For the pair inequalities, define
\begin{align*}
 D_0(p,r)&=\frac{A(p)A(r)(A(p)+A(r))}{1000},\\
 K_a&=24\bigl(A(p)+A(r)-H_4(p)F_4(r)-F_4(p)H_4(r)-D_0(p,r)\bigr),\\
 K_b&=-24\bigl(H_4(p)G_4(r)+G_4(p)H_4(r)\bigr).
\end{align*}
The two targets have the common form
\begin{equation}
 \mathcal R=K+\sum_{\pi\in S_4}\frac{N_\pi}{z_\pi^k},
 \label{eq:q4-rational-target}
\end{equation}
with the following data:
\begin{center}
\begin{tabular}{cclc}
\hline
Target & $k$ & $N_\pi$ & $m$\\
\hline
$a$ & 1 & $-\lVert LU^\pi\rVert_2^2$ & $3$\\
$b$ & 5 & $z_\pi^6H_4(U^\pi/z_\pi)$ & $1/2$\\
\hline
\end{tabular}
\end{center}
The numerator in the second row is polynomial: each degree-$s$ term of
$H_4$ is multiplied by $z_\pi^{6-s}$ after substitution.  The two values of
$\mathcal R$ are exactly twenty-four times the respective nonnegative
differences in \eqref{eq:q4-pair-A} and \eqref{eq:q4-pair-H}.

The verifier covers $\mathcal S_4\times\mathcal S_4$ by products of rational
subsimplices.  On each cell it applies one of two sufficient, exact
coefficient tests.

The first test uses
\[
 t_k(z)=\sum_{j=0}^5\binom{-k}{j}(z-1)^j\le z^{-k}
 \qquad(z>0).
\]
Taylor's theorem proves the inequality on the entire positive half-line:
the sixth derivative of $z^{-k}$ is positive, and the remainder contains
the nonnegative factor $(z-1)^6$.  The scalar bounds above give
$N_\pi+m z_\pi^{k+1}\ge0$.  Hence
\[
 \frac{N_\pi}{z_\pi^k}
 \ge (N_\pi+m z_\pi^{k+1})t_k(z_\pi)-m z_\pi.
\]
Nonnegative homogeneous coefficients of the sum of these polynomial lower
bounds and $K$ certify the cell.  The nonnegative multiplier is necessary;
one cannot multiply a lower bound by an uncontrolled signed numerator.

For the second test, choose rational bihomogeneous polynomials $J_\pi$ of
bidegree $(3,3)$.  With
$W_j(X,Y)=(\sum_iX_i)^j(\sum_iY_i)^j$, form
\[
 O_\pi=W_2N_\pi-z_\pi^kJ_\pi,
 \qquad Q_\pi=z_\pi^kW_3.
\]
These are both of bidegree $(k+3,k+3)$, and $Q_\pi$ has nonnegative
coefficients.  The checker chooses a rational $c_\pi\ge0$ such that every
coefficient of $O_\pi+c_\pi Q_\pi$ is nonnegative.  A negative coefficient
of $O_\pi$ at a zero coefficient of $Q_\pi$ makes this test fail.  Otherwise
the required correction is obtained by rounding each necessary coefficient
ratio upward to a multiple of $2^{-32}$ and taking the largest correction.
On the interior of the cell this proves
\[
 \frac{N_\pi}{z_\pi^k}
 =J_\pi-c_\pi+
   \frac{O_\pi+c_\pi Q_\pi}{z_\pi^k}
 \ge J_\pi-c_\pi.
\]
The cell is accepted when the homogeneous coefficients of
$K+\sum_\pi(J_\pi-c_\pi)$ are all nonnegative, after degree elevation.
The polynomials $J_\pi$ in the source are obtained by rational interpolation
on the degree-three simplex lattice.  Their interpolation accuracy is not
assumed: the full remainder is certified by the displayed identity.

Unaccepted cells are bisected at an edge midpoint.  If vertex $i$ is
replaced by the midpoint of vertices $i,j$, the barycentric substitution is
\[
 X_i=X_i'/2,\qquad X_j=X_j'+X_i'/2.
\]
The binomial expansion gives exact rational coefficients for the child
cell.  The two choices of the replaced endpoint cover the original
simplex, corresponding to $X_i\le X_j$ and $X_j\le X_i$.  Repeating this
construction therefore proves coverage, rather than merely testing a set
of posterior points.

The full integer verifier of \cite{OpenAI4}, with assertions enabled, was
replayed for both targets.  It terminated with the following counts and no
unresolved cells:
\begin{center}
\begin{tabular}{crrrr}
\hline
Target & Visited cells & Maximum depth & Taylor leaves & Remainder leaves\\
\hline
$a$ & 4,479 & 19 & 166 & 2,074\\
$b$ & 1,935 & 20 & 97 & 871\\
\hline
\end{tabular}
\end{center}
The scalar coefficients and the twenty-four induced coordinate actions were
also checked independently by rational arithmetic.  All sign decisions in
the certificates use integer numerators and positive exact denominators.

Finally, denominator clearing and the lower bounds above are used first on
cell interiors, where every $z_\pi>0$.  The weighted terms
$z_\pi A(P^\pi)$ and $z_\pi H_4(P^\pi)$ converge to zero when
$z_\pi\to0$, since $A,H_4$ are bounded on $\Delta_4$.  They therefore have
continuous extensions to every boundary point.  This proves the boundary
cases without dividing by a vanishing denominator.  The original source
contains the complete verifier and its deterministic subdivision rule.

\subsection{Closure of the four-state cone}

Let $\mathcal C_4$ consist of the permutation-symmetric posterior laws
$\mu$ satisfying $\mathbb E_\mu H_4\ge0$.  The condition is imposed on the
law, not on each posterior in its support.  Full revelation belongs to the
cone because $H_4(e_i)=13/100$; the uninformative experiment belongs because
$H_4(u_4)=0$.

\begin{lemma}[Four-state experiment closure]
\label{lem:q4-cone-closure}
The class $\mathcal C_4$ is preserved by a channel edge with
$0\le\lambda\le1$, by conditionally independent products, and by observed
mixtures with an input-independent mixing label.  If
$\mu,\nu\in\mathcal C_4$, and
$a=\mathbb E_\mu A$, $b=\mathbb E_\nu A$, then
\begin{equation}
 \mathbb E_{\mu\star\nu}A
 \le a+b-
 \frac{(\mathbb E_\mu A^2)b+a(\mathbb E_\nu A^2)}{1000}
 \le a+b-\frac{ab(a+b)}{1000}.
 \label{eq:q4-experiment-saving}
\end{equation}
\end{lemma}

\begin{proof}
An edge sends $p$ to $u_4+\lambda(p-u_4)$, and hence sends $Lp$ to
$\lambda Lp$.  Equation~\eqref{eq:q4-radial} preserves the expectation
condition, and $A$ is multiplied by $\lambda^2$.

For the product of two experiments, the expectation of a posterior
functional $f$ is
\[
 \mathbb E_{\mu\star\nu}f
 =\mathbb E_{\mu\otimes\nu}[z f(P)].
\]
The laws in the reference measure $\mu\otimes\nu$ are independent; the
factor $z$ gives the actual joint observation law.  By permutation
symmetry, averaging the integrand over the relative permutations does not
change the expectation.  Integrating \eqref{eq:q4-pair-H} therefore gives
\[
 \mathbb E_{\mu\star\nu}H_4
 \ge (\mathbb E_\mu H_4)(\mathbb E_\nu G_4)
      +(\mathbb E_\mu G_4)(\mathbb E_\nu H_4)\ge0.
\]
The product is symmetric as well.  Integrating
\eqref{eq:q4-pair-A} and discarding the nonnegative products involving
$F_4$ gives the first inequality in
\eqref{eq:q4-experiment-saving}.  Jensen's inequality gives the second.
Observed input-independent mixtures preserve both symmetry and the
expectation condition by linearity.
\end{proof}

\section{Nonreconstruction for arbitrary finite-mean offspring}
\label{sec:finite-mean-nonreconstruction}

Let $D$ be a nonnegative, almost surely finite integer-valued random
variable with $d=\mathbb E D<\infty$, and let $T$ be the associated
Galton--Watson tree.  The tree and all auxiliary child orderings are
independent of the spins.  The observer knows the tree and the spins at
generation $n$.  Knowing the continuation of the tree beyond that
generation adds no spin information conditional on the observed finite
prefix.

For $q=3,4$, write $p_n$ for the posterior of the uniformly distributed
root spin.  Use the quadratic statistics
\[
 J_3(p)=\frac{3\sum_{i=1}^3p_i^2-1}{2},\qquad
 J_4(p)=4\sum_{i=1}^4p_i^2-1,
 \qquad m_n^{(q)}=\mathbb E J_q(p_n).
\]
Thus $J_3=x$ in the three-state section and $J_4=A$ in the preceding
section.  Initially $m_0^{(3)}=1$ and $m_0^{(4)}=3$.  Set
\begin{equation}
 \rho=d\lambda^2,\qquad
 \delta=\mathbb E(D-1)_+=d-\mathbb P(D\ge1).
 \label{eq:gw-rho-delta}
\end{equation}
Only the first offspring moment is involved in $\delta$.

\begin{proposition}[Quadratic and cubic drift]
\label{prop:gw-drift}
For $0\le\lambda\le1$ and every $n\ge0$,
\begin{align}
 m_{n+1}^{(3)}
 &\le \rho m_n^{(3)}-c_3\bigl(m_n^{(3)}\bigr)^2,
 &c_3&=\frac{\delta\lambda^4}{1000},
 \label{eq:gw-q3-drift}\\
 m_{n+1}^{(4)}
 &\le \rho m_n^{(4)}-c_4\bigl(m_n^{(4)}\bigr)^3,
 &c_4&=\frac{\delta\lambda^6}{500}.
 \label{eq:gw-q4-drift}
\end{align}
\end{proposition}

\begin{proof}
At depth zero the fully revealing experiment belongs to the appropriate
cone.  Edge attenuation, conditionally independent products, and observed
offspring mixtures show inductively that every finite-depth experiment
belongs to that cone.  These closures also apply to every partial product
of incoming branch experiments.

Condition on $D=k$ and put $b=\lambda^2m_n^{(q)}$.  Every child branch
has mean statistic $b$, since one channel edge multiplies the quadratic
statistic by $\lambda^2$.  If $j\ge1$ branches have already been observed,
their partial product has mean statistic at least $b$.  To see this,
let $P^{(j)}$ be the posterior after those $j$ branches and $P^{(1)}$ the
posterior after the first branch.  The tower property gives
\[
 P^{(1)}=\mathbb E[P^{(j)}\mid\hbox{first branch observation}],
\]
and conditional Jensen applies because both $J_3$ and $J_4$ are convex.

The three-state pair inequality proved earlier gives, for a partial
product with mean $a\ge b$ and one further branch,
\[
 m(\hbox{new product})\le a+b-\frac{ab}{1000}
 \le a+b-\frac{b^2}{1000}.
\]
Applying this at each of the $k-1$ fusions, and including $k=0,1$, yields
\[
 \mathbb E[J_3(p_{n+1})\mid D=k]
 \le kb-\frac{(k-1)_+b^2}{1000}.
\]
For four states, \eqref{eq:q4-experiment-saving} instead gives a saving
of at least
\[
 \frac{ab(a+b)}{1000}\ge\frac{b^3}{500}
\]
at each fusion.  Therefore
\[
 \mathbb E[J_4(p_{n+1})\mid D=k]
 \le kb-\frac{(k-1)_+b^3}{500}.
\]
Averaging over the observed offspring count proves both recursions.
The expectation is finite because $\mathbb E D<\infty$.
\end{proof}

\begin{proposition}[Decay at and below the Kesten--Stigum value]
\label{prop:gw-decay}
Assume $0<\lambda<1$ and $\rho\le1$.  Let
\[
 a_n^{(q)}=\mathbb E\operatorname{TV}(p_n,u_q).
\]
If $\rho<1$, then
\begin{align}
 m_n^{(3)}&\le\rho^n,
 &a_n^{(3)}&\le\frac{\rho^{n/2}}{\sqrt2},
 \label{eq:q3-subcritical-rate}\\
 m_n^{(4)}&\le3\rho^n,
 &a_n^{(4)}&\le\frac{\sqrt3}{2}\rho^{n/2}.
 \label{eq:q4-subcritical-rate}
\end{align}
If $\rho=1$, then $c_3,c_4>0$ and
\begin{align}
 m_n^{(3)}&\le\frac1{1+c_3n},
 &a_n^{(3)}&\le\frac1{\sqrt{2(1+c_3n)}},
 \label{eq:q3-critical-rate}\\
 m_n^{(4)}&\le(1/9+2c_4n)^{-1/2},
 &a_n^{(4)}&\le\frac12(1/9+2c_4n)^{-1/4}.
 \label{eq:q4-critical-rate}
\end{align}
In both cases nonreconstruction holds for almost every realized tree.
\end{proposition}

\begin{proof}
Below criticality, omit the negative terms in
\eqref{eq:gw-q3-drift}--\eqref{eq:gw-q4-drift} and iterate.  The
total-variation estimates follow from Cauchy--Schwarz:
\[
 \operatorname{TV}(p,u_3)\le\sqrt{J_3(p)/2},\qquad
 \operatorname{TV}(p,u_4)\le\frac12\sqrt{J_4(p)},
\]
followed by Jensen's inequality.

At $\rho=1$, the condition $\lambda<1$ implies $d>1$, and hence
$\delta=d-\mathbb P(D\ge1)>0$.  For three states, while
$m_n=m_n^{(3)}>0$,
\[
 \frac1{m_{n+1}}
 \ge\frac1{m_n(1-c_3m_n)}
 \ge\frac1{m_n}+c_3.
\]
The denominator is positive: $m_n\le1$ and
$c_3\le d\lambda^4/1000\le1/1000$ when $\rho\le1$.
If an iterate is zero, all subsequent iterates are zero and the stated
bound is immediate.  Otherwise iteration from $m_0=1$ proves
\eqref{eq:q3-critical-rate}.

For four states, write $m_n=m_n^{(4)}$.  Since $m_n\le3$ and
$c_4m_n^2\le9d\lambda^6/500\le9/500<1$, the analogous calculation gives
\[
 \frac1{m_{n+1}^2}
 \ge\frac1{m_n^2}(1-c_4m_n^2)^{-2}
 \ge\frac1{m_n^2}+2c_4.
\]
Iteration from $m_0=3$, with the same zero-iterate convention, proves
\eqref{eq:q4-critical-rate}.

For a fixed locally finite tree $T$, put
\[
 a_n^{(q)}(T)=
 \mathbb E[\operatorname{TV}(p_n,u_q)\mid T].
\]
The root spin, level $n$, and level $n+1$ form a Markov chain, so
$a_n^{(q)}(T)$ is nonincreasing by data processing.  Its nonnegative limit
therefore exists.  The annealed estimates show that its expectation tends
to zero.  Bounded convergence yields
\[
 \mathbb{E}\bigl[\lim_{n\to\infty}a_n^{(q)}(T)\bigr]=0,
\]
so the limit is zero for almost every tree.  The conclusion also holds
almost surely conditional on survival whenever survival has positive
probability.  No independence assumption after conditioning on survival
is needed.
\end{proof}

\section{Reconstruction above the threshold}
\label{sec:finite-mean-supercritical}

The argument in this section works for every finite number of states
$q\ge2$.  It supplies the supercritical direction with only a finite first
offspring moment.

\begin{proposition}[Bounded pruning]
\label{prop:bounded-pruning}
Let $\mathbb E D=d<\infty$ and $0<\lambda<1$.  If
$d\lambda^2>1$, then the annealed reconstruction advantage is bounded
away from zero uniformly in depth.
\end{proposition}

\begin{proof}
By monotone convergence, choose $L<\infty$ such that
\[
 D_L=\min(D,L),\qquad \mu=\mathbb E D_L,
 \qquad \mu\lambda^2>1.
\]
Keep the first $D_L$ children at every retained vertex, using a child
ordering independent of the spins.  The retained subtree has iid bounded
offspring $D_L$, and its depth-$n$ spins are part of the original
depth-$n$ observation.

Choose a channel eigenfunction $\phi$ with uniform mean zero and second
moment one; for example,
\[
 \phi(1)=\sqrt{q-1},\qquad
 \phi(i)=-\frac1{\sqrt{q-1}}\quad(2\le i\le q).
\]
It satisfies
$\mathbb E[\phi(\sigma_{\rm child})\mid\sigma_{\rm parent}]
=\lambda\phi(\sigma_{\rm parent})$ and
$|\phi|\le B_q:=\sqrt{q-1}$.
On the retained tree define
\[
 Z_n=(\mu\lambda)^{-n}
       \sum_{|v|=n}\phi(\sigma_v).
\]
The first-moment recursion gives
\begin{equation}
 \mathbb E[Z_n\mid\sigma_{\rm root}]
 =\phi(\sigma_{\rm root}),\qquad
 \mathbb E[\phi(\sigma_{\rm root})Z_n]=1.
 \label{eq:pruning-first-moment}
\end{equation}

Let $M_n=\mathbb E Z_n^2$ and
$\alpha_L=\mathbb E[D_L(D_L-1)]/\mu^2$.  Writing the depth-$(n+1)$
sum as a sum of child contributions gives
\begin{equation}
 M_{n+1}=\frac{M_n}{\mu\lambda^2}+\alpha_L,
 \qquad M_0=1.
 \label{eq:pruning-second-moment}
\end{equation}
Indeed the diagonal terms have total expectation $\mu M_n$.  For distinct
children, conditional independence given the root and
\eqref{eq:pruning-first-moment} give product expectation
$\lambda^2\phi(\sigma_{\rm root})^2$, whose mean is $\lambda^2$.
Dividing by $(\mu\lambda)^2$ proves
\eqref{eq:pruning-second-moment}.  Since $D_L$ is bounded and
$r_L=(\mu\lambda^2)^{-1}<1$,
\[
 M_n\le C_L:=1+\frac{\alpha_L}{1-r_L}<\infty
 \qquad\hbox{for all }n.
\]

Let $\mathcal F_n$ contain the original tree, its level-$n$ spins, and
the auxiliary ordering.  The statistic $Z_n$ is $\mathcal F_n$-measurable.
Put $\eta_n=\mathbb E[\phi(\sigma_{\rm root})\mid\mathcal F_n]$.
Cauchy--Schwarz and \eqref{eq:pruning-first-moment} imply
\[
 1=\mathbb E[\eta_nZ_n]
 \le\bigl(\mathbb E\eta_n^2\,C_L\bigr)^{1/2},
 \qquad \mathbb E\eta_n^2\ge C_L^{-1}.
\]
The ordering is independent of the root and spins conditional on the
original observed tree.  Consequently adjoining it does not change the
root posterior.  For that posterior,
\[
 |\eta_n|\le B_q,\qquad
 \operatorname{TV}(p_n,u_q)\ge\frac{|\eta_n|}{2B_q}.
\]
Therefore
\begin{equation}
 a_n^{(q)}
 \ge\frac{\mathbb E|\eta_n|}{2B_q}
 \ge\frac{\mathbb E\eta_n^2}{2B_q^2}
 \ge\frac1{2(q-1)C_L}>0.
 \label{eq:supercritical-advantage}
\end{equation}
Only the truncated offspring variable was required to have a second
moment.
\end{proof}

\begin{lemma}[The reconstruction event on a Galton--Watson tree]
\label{lem:survival-reconstruction}
For $0<\lambda<1$ and $d\lambda^2>1$, reconstruction holds for almost
every surviving Galton--Watson tree and fails on the extinction event.
\end{lemma}

\begin{proof}
First fix a deterministic locally finite rooted tree.  Its root
reconstructs if and only if at least one of its child subtrees reconstructs.
If a child reconstructs, observing only that child's descendants gives
the root posterior
\[
 u_q+\lambda(p_{\rm child}-u_q).
\]
Its total-variation advantage is $\lambda$ times the child's advantage,
so remains positive in the limit.

Conversely, suppose all child subtrees are nonreconstructing.  For a finite
input alphabet, vanishing expected posterior total variation implies
vanishing mutual information: the posterior relative entropy from the
uniform prior is bounded and continuous on the compact simplex.  Thus
every child's boundary observation has vanishing mutual information with
the child spin and, by data processing, with the root spin.  There are
only finitely many children.  Their observations $Y_1,\ldots,Y_k$ are
independent conditional on the root, and hence
\[
 I(\sigma_{\rm root};Y_1,\ldots,Y_k)
 \le\sum_{j=1}^k I(\sigma_{\rm root};Y_j)\longrightarrow0.
\]
Here the inequality follows from subadditivity of unconditional entropy
and additivity of conditional entropy.  Pinsker's inequality then implies
nonreconstruction at the root.

Let $r_* $ be the probability that a Galton--Watson tree reconstructs.
The preceding equivalence and independence of child subtrees yield
\[
 r_*=1-\mathbb E(1-r_*)^D.
\]
Let $\xi$ be the extinction probability.  Since
$d\lambda^2>1$ implies $d>1$, the offspring generating function has
exactly the two fixed points $\xi$ and $1$ in $[0,1]$.  For completeness,
it is strictly convex on $(0,1)$ because $\mathbb P(D\ge2)>0$;
its derivative at $1$ equals $d>1$, and $\xi$ is its least fixed point.
These facts exclude any third fixed point.

The positive lower bound \eqref{eq:supercritical-advantage}, together
with bounded convergence for the fixed-tree advantage, implies $r_*>0$.
Thus $1-r_*=\xi$.  Reconstruction is impossible on an extinct tree,
whose boundary is eventually empty.  Consequently the reconstruction
event agrees with survival up to a null set.
\end{proof}

\section{The finite-mean dichotomy and endpoint conventions}
\label{sec:full-gw-dichotomy}

\begin{theorem}[Three and four states on finite-mean Galton--Watson trees]
\label{thm:full-finite-mean-dichotomy}
Let $D$ be a nonnegative, almost surely finite integer-valued random
variable with finite mean $d$.  Consider the symmetric ferromagnetic
$q$-state broadcast process on its observed Galton--Watson tree, with
$q\in\{3,4\}$ and $0<\lambda<1$.
If $d\lambda^2\le1$, nonreconstruction holds for almost every tree.
If $d\lambda^2>1$, reconstruction holds for almost every surviving
tree and fails on extinct trees.  In particular, annealed reconstruction
occurs if and only if
\[
 d\lambda^2>1.
\]
No offspring moment beyond $\mathbb E D<\infty$ is assumed.  At equality,
the quantitative bounds are \eqref{eq:q3-critical-rate} and
\eqref{eq:q4-critical-rate}.
\end{theorem}

\begin{proof}
Proposition~\ref{prop:gw-decay} proves the assertion at and below equality.
Proposition~\ref{prop:bounded-pruning} and
Lemma~\ref{lem:survival-reconstruction} prove the assertions above equality.
\end{proof}

\begin{remark}[Regular trees and endpoints]
Taking $D\equiv d$ gives every regular tree of integer degree $d\ge2$.
At $\lambda=0$, all positive-depth observations are independent of the
root.  At $\lambda=1$, every descendant copies the root, so reconstruction
is perfect on survival and impossible after extinction.  A subcritical
or nondegenerate critical finite-mean Galton--Watson process dies out
almost surely.  The exceptional law $D\equiv1$ instead gives an infinite
noiseless line: at $\lambda=1$ it reconstructs perfectly although
$d\lambda^2=1$.  A statement including all closed endpoints must record
this exception.  The argument concerns nonnegative channel eigenvalues;
it makes no assertion about the antiferromagnetic case.
\end{remark}
